% Part: lambda-calculus
% Chapter: introduction
% Section: lambda-computable

\documentclass[../../../include/open-logic-section]{subfiles}

\begin{document}

\olfileid{lam}{int}{cmp}
\olsection{\usetoken{S}{lambda definable} વિધેયો સંગણનીય છે}

\begin{thm}
\ollabel{thm:lambda-computable}
આંશિક વિધેય~$f$ લૅમ્બડા પદ વડે !!{lambda defined} હોય, તો તે
સંગણનીય છે.
\end{thm}

\begin{proof}
માનો કે વિધેય~$f$ લૅમ્બડા પદ~$X$ વડે !!{lambda defined} છે.
$f$ ગણવાની અનૌપચારિક પ્રક્રિયા વર્ણવીએ. ઇનપુટ $m_0$,
\dots,~$m_{n-1}$ માટે પદ $X \num{m_0} \ldots \num{m_{n-1}}$
લખો. ત્યાર પછી વૃક્ષ બનાવો: પહેલાં મૂળ પદનાં બધાં એક-પગલાનાં
લઘુકરણો લખો; તેમની નીચે તે પદોનાં બધાં એક-પગલાનાં લઘુકરણો
(અર્થાત્ મૂળ પદનાં બે-પગલાનાં લઘુકરણો) લખો; અને એમ આગળ વધો.
ક્યારેય ચર્ચ અંક મળે તો તેને ઉત્તર તરીકે પરત કરો; ન મળે તો
વિધેયનું મૂલ્ય અવ્યાખ્યાયિત રહે છે.
\sourcecorrection{OLLAM-005}{સ્થિર અંગ્રેજી મૂળમાં પ્રથમ
ઇનપુટ-પદમાં~\texttt{\textbackslash num m\_0} અને
\texttt{\textbackslash num m\_\{n-1\}} છે. આથી~$m_0$
અને~$m_{n-1}$ આખાં \texttt{\textbackslash num}ના દલીલરૂપે
જતા નથી. નીચેની જ પંક્તિમાં વપરાતી સાચી શૈલી મુજબ અહીં
\texttt{\textbackslash num\{m\_0\}} વગેરે લખ્યું છે.}

ચર્ચની માન્યતાનો આધાર લઈને કહી શકાય કે આ વિધેય સંગણનીય છે.
પ્રમેયની વધુ સારી સાબિતી માટે આ શોધપ્રક્રિયાનું પુનરાવર્તી વર્ણન
આપવું જોઈએ. દાખલા તરીકે,
``$\fn{IsASubterm}$'', ``$\fn{Substitute}$'',
``$\fn{ReducesToInOneStep}$'', ``$\fn{ReductionSequence}$'',
``$\fn{Numeral}$'' વગેરે આદિમ પુનરાવર્તી વિધેયો અને સંબંધોની
શ્રેણી વ્યાખ્યાયિત કરી શકાય. ત્યાર પછી $f(m_0, \dots, m_{n-1})$
ગણવાની આંશિક પુનરાવર્તી પ્રક્રિયા એ છે કે
$X \num{m_0} \dots \num{m_{n-1}}$થી શરૂ થઈ ચર્ચ અંક પર પૂરી થતી
એક-પગલાનાં લઘુકરણોની શ્રેણી શોધો, અને તે અંકને અનુરૂપ સંખ્યા
પરત કરો. વિગત લાંબી અને કંટાળાજનક છે, પરંતુ બાકી નિયમિત છે.
\end{proof}

\end{document}
